% Источники: exam/materials/моделирование.pdf, раздел 2; exam/materials/прошлогодние ответы/Ответы_Моделирование.pdf, вопрос 2; https://github.com/HumsterProgrammer/iu9-notes/blob/main/8sem/Моделирование/Моделирование_лекция-01_09-02-26.md; https://github.com/HumsterProgrammer/iu9-notes/blob/main/8sem/Моделирование/Моделирование_лекция-02_12-02-26.md.
\section{Решение инженерных задач на компьютере.}

Решение инженерных задач на компьютере является важной частью современного инженерного дела. Инженерные задачи, как правило, имеют ярко выраженную техническую направленность: разработка техпроцесса, создание конструкции, оптимизация затрат на производстве.

Основные характеристики инженерных задач:
\begin{enumerate}
  \item инженерные задачи имеют ярко-выраженную практическую направленность (создание новых конструкций, разработка техпроцесса, минимизация затрат на производство изделий и т.д.);
  \item для инженерных задач характерна необходимость доведения результатов до конкретных чисел, на основании которых можно принимать решения;
  \item для инженерных задач характерен значительный объем выполняемой вычислительной работы, при этом набор методов и их реализаций конечен;
\end{enumerate}

В основе решения инженерных задач лежит теория инженерного эксперимента, включающая три раздела: планирование эксперимента, проведение эксперимента, обработка результатов эксперимента.

Этапы решения инженерной задачи на компьютере:
\begin{enumerate}
  \item постановка задачи: определение входных данных, требований к ним и к математической модели;
  \item выбор или построение математической модели с учетом компромисса между полнотой описания и сложностью реализации;
  \item декомпозиция математической модели: разбиение на набор вычислительных задач;
  \item предмашинный анализ свойств вычислительной задачи: оценка корректности, обусловленности, устойчивости;
  \item выбор и построение численного метода или методов для решения выделенных вычислительных задач;
  \item алгоритмизация и программирование численного метода или методов;
  \item отладка программы;
  \item тестирование программы;
  \item обработка и интерпретация результатов;
  \item использование результатов моделирования или возврат к построению/модификации модели.
\end{enumerate}

\clearpage
